% T11 — «Тетрадь»: линованный фон, «рукописный» заголовок (sans bold italic), синий.
% Профиль: математика, неравенства. 6 задач, 1 кол.
\documentclass[a4paper,11pt]{article}

\usepackage{../../../../Lessons/_templates/preamble_worksheet}
\usepackage{../_shared/worksheet_check}

\usepackage{eso-pic}  % фоновый рисунок на каждой странице

% --- Палитра: чернильный синий ---
\definecolor{wcprimary}{HTML}{2C5282}
\definecolor{wcprimaryLight}{HTML}{BEE3F8}
\definecolor{wcprimaryBG}{HTML}{EBF8FF}
\definecolor{notebookline}{HTML}{C3DAFE}
\definecolor{notebookmargin}{HTML}{FC8181}

% --- Фон: горизонтальные линии тетради + вертикальная красная полоса слева ---
\AddToShipoutPictureBG{%
  \begin{tikzpicture}[remember picture, overlay]
    % Горизонтальные линии каждые 7 мм
    \foreach \y in {1,...,40}
      \draw[notebookline, line width=0.3pt]
        (current page.north west) ++(20mm, -\y*7mm)
        -- ([xshift=-12mm, yshift=-\y*7mm]current page.north east);
    % Красная вертикальная «поля»
    \draw[notebookmargin, line width=0.6pt]
      ([xshift=18mm]current page.north west) -- ([xshift=18mm]current page.south west);
    \draw[notebookmargin, line width=0.6pt]
      ([xshift=20mm]current page.north west) -- ([xshift=20mm]current page.south west);
  \end{tikzpicture}}

\geometry{a4paper, top=18mm, bottom=18mm, left=26mm, right=14mm}

% --- Заголовок: «рукописный» (sans + italic + bold) ---
\renewcommand{\WorksheetTitle}[1]{%
  \par\noindent
  {\sffamily\itshape\bfseries\fontsize{22pt}{26pt}\selectfont\color{wcprimary}#1}\par
  \vspace{2mm}{\color{wcprimary!60}\rule{0.5\linewidth}{0.6pt}}\par\vspace{3mm}}

% --- TaskBox: только «№N.» жирным, без рамки, текст поверх линий ---
\renewcommand{\TaskBox}[2]{%
  \par\nopagebreak\noindent
  {\sffamily\bfseries\color{wcprimary}№#1.}\ #2\par
  \vspace{6mm}}

\SetAnswerFieldSize{30mm}{9mm}

\begin{document}

\WorksheetTitle{Линейные неравенства}
{\sffamily\itshape\small\color{wcprimary!80}Самостоятельная работа \quad·\quad Класс 8 \quad·\quad T11 \quad·\quad ID: T11-ineq-2026-05-26}\par\vspace{3mm}
\FirstPageHeader

\TaskBox{1}{Решите неравенство $3x - 7 > 2x + 5$. В ответе запишите наименьшее целое решение.\par\vspace{1mm}\hfill\textbf{Ответ:}~\AnswerField{1}}

\TaskBox{2}{Решите неравенство $5(x - 1) \leqslant 2(2x + 3)$. В ответе запишите наибольшее целое решение.\par\vspace{1mm}\hfill\textbf{Ответ:}~\AnswerField{2}}

\TaskBox{3}{Решите систему $\begin{cases} 2x - 5 > 1 \\ 7 - x > 0 \end{cases}$. Сколько целых решений?\par\vspace{1mm}\hfill\textbf{Ответ:}~\AnswerField{3}}

\TaskBox{4}{Решите неравенство $\dfrac{x - 4}{3} \geqslant \dfrac{2x + 1}{6}$. В ответе запишите наибольшее целое решение.\par\vspace{1mm}\hfill\textbf{Ответ:}~\AnswerField{4}}

\TaskBox{5}{При каких значениях $a$ неравенство $ax > 5$ имеет решение $x < 2$? Если $a$ единственное, запишите его; иначе $0$.\par\vspace{1mm}\hfill\textbf{Ответ:}~\AnswerField{5}}

\TaskBox{6}{Найдите количество целых решений двойного неравенства $-3 \leqslant 2x - 1 < 7$.\par\vspace{1mm}\hfill\textbf{Ответ:}~\AnswerField{6}}

\WorksheetQR{T11-ineq/L1/v1}

\end{document}
